\section{Discussion}\label{sec:discussion}

Our first result replicates, and amplifies, the central finding of \citet{knaap2024segregated}. When local
environments are defined by walking distance, racial segregation in large Brazilian cities is higher in every
case, by \meanPdH\% on average, and the ranking of the most segregated places changes. Cities built on hills,
islands and valleys (Florian\'opolis, Vit\'oria, Juiz de Fora) climb the ranking. Analysts who compare cities
with Euclidean spatial indices would therefore reach different conclusions about where segregation is most
intense. This is consistent with the Brazilian configurational literature, which describes Brazilian street
networks as unusually fragmented and labyrinthine \citep{medeiros2013urbis,holanda2002espaco}.

Our second result qualifies the first. Once the population of local environments is held constant, the gap
vanishes. Walking along the network does not reveal a different social composition next door. It reveals a
smaller neighbourhood, and smaller neighbourhoods are, as the multiscalar literature shows, more segregated
\citep{reardon2008geographic}. This supports the concern raised by \citet{saporito2026reconsidering} and changes
how the network--Euclidean gap should be read. It is not evidence that streets act as barriers that sort groups
within a given population radius. It is evidence that street networks set the \emph{effective scale} at which
residents experience their city: in dendritic, circuitous, cul-de-sac-rich fabrics, a 2-km walk reaches
as little as \ratioMin\% of the people a 2-km circle would. That is still a finding about urban design. Where the network
is poorly connected, the local environment of an average resident is smaller and therefore more homogeneous,
which matters for everyday exposure to other groups \citep{netto2015segregated}. But it calls for a different
framing, and for population-based (or population-matched) comparisons whenever network and Euclidean indices are
contrasted.

Third, we find no support for configurational fragmentation, as we operationalise it, as a predictor of the gap.
Neither the gap nor the share of reachable population is associated with $F$. What matters is the local
texture of the network: dead ends, circuity and meshedness. These are exactly the features that gated
condominiums and closed \emph{loteamentos} produce \citep{coy2006gated}. One interpretation is that
local-environment indices at a 2-km scale are sensitive to local permeability and blind to city-wide
articulation, the property that synergy captures. Another is that our metric analogue does not capture what
space-syntax synergy measures on axial maps. Testing CHASM-type segment-based measures \citep{saboya2025chasm},
which build the street grid into exposure itself, is a natural next step.

\paragraph{Limitations.} First, OpenStreetMap coverage of informal settlements is uneven: stairways and alleys
in favelas are often missing, which lengthens network paths and may inflate the gap there, so our estimates for
hillside cities are best read as upper bounds. Second, tracts are represented by a single interior point, and
\nSnapFar{} of the \nTractsUrban{} tracts lie more than 500~m from the nearest network node. Third, we consider the
walking network only; public transport would lengthen reach and probably reduce the gap
\citep{pereira2021r5r,bittencourt2021cumulative}. Fourth, the units are municipalities, not functional urban
areas, and results are subject to the modifiable areal unit problem. Finally, following the design of this paper,
we do not test whether the gaps are larger than would be expected by chance. The random-labelling framework of
\citet{rey2021comparative} and \citet{knaap2024segregated} can be applied directly to our cached weights.

\section{Conclusion}\label{sec:conclusion}

Measured with walking distance, racial segregation in the \nCities{} largest Brazilian municipalities is
uniformly higher than when measured as the crow flies, and rankings of the most segregated cities change. But
the difference disappears once local environments are matched in population. The street network shapes measured
segregation mainly by shrinking the neighbourhood rather than by separating groups within it. Dead ends and
circuity, not city-wide configurational fragmentation, are associated with that shrinkage. For measurement, this
means that network and Euclidean indices should be compared at equal population scale. For planning, it points to
the permeability of the local street fabric, the connectivity of peripheral \emph{loteamentos} and gated enclaves,
as the lever through which urban design conditions everyday exposure between racial groups. Future work should
add inferential tests of the gap, multiscalar profiles, and public-transport networks.
